\documentclass[11pt]{article}

% --------- PACKAGES ----------
\usepackage[left=0.8in,top=0.7in,right=0.8in,bottom=0.7in]{geometry}
\usepackage{enumitem}
\usepackage[hidelinks]{hyperref}
\usepackage{setspace}
\usepackage{parskip}
\usepackage[en-GB]{datetime2}
\usepackage[T1]{fontenc}
\usepackage{newtxtext,newtxmath}   % Times-like
% pdfTeX-only: improves ATS text extraction. Guarded so XeTeX/Tectonic skip it.
\ifdefined\pdfgentounicode
  \input{glyphtounicode}
  \pdfgentounicode=1
\fi

\setlength{\parindent}{0pt}
\setstretch{1.05}
\pagenumbering{gobble}

% --------- FIXED IDENTITY (KG §1 — do not edit) ----------
\newcommand{\clName}{Nishal Stanislaus Silva}
\newcommand{\clEmail}{nishal.silva@hotmail.com}
\newcommand{\clPhone}{+1 613 200 2030}
\newcommand{\clCity}{Toronto, ON}
\newcommand{\clSite}{https://nishal.xyz}
\newcommand{\clLinkedIn}{https://www.linkedin.com/in/nishal-silva}
\newcommand{\clStatus}{Canadian Permanent Resident, Eligible to work in Canada without sponsorship}

% --------- PER-APPLICATION VARIABLES ----------
\newcommand{\clTagline}{Machine Learning Research Engineer | Computer Vision, Deep Learning, Real-Time Edge Inference | Applied ML}
\newcommand{\clRole}{AI Vision Engineer}
\newcommand{\clCompany}{Nexxa.AI}
\newcommand{\clAddressee}{Dear Hiring Manager,}

\begin{document}
\pagestyle{empty}

\begin{center}
    {\LARGE \scshape \textbf{\clName}}{\large \textit{, PhD}}
    \\[1pt]
    {\small \clTagline}
    \\[1pt]
    {\footnotesize\textit{\clStatus}}
    \\[1pt]
    {\small
    \href{mailto:\clEmail}{\clEmail}
    \,\,|\,\, \clPhone
    \,\,|\,\, \clCity
    \,\,|\,\, \href{\clSite}{nishal.xyz}
    \,\,|\,\, \href{\clLinkedIn}{linkedin.com/in/nishal-silva}}
\end{center}

\rule{\linewidth}{0.4pt}\vspace{0.6em}

\today

\textbf{Re: \clRole{} -- \clCompany{}}

\clAddressee

I am applying for the AI Vision Engineer role at Nexxa.AI. I am a machine learning engineer with a PhD from the University of Trento and 10+ years across industry and academia. For five and a half years at MAS Holdings I led production computer vision for manufacturing, from care-label inspection to loom-side fabric defect detection, and that work is the closest match to what you describe: vision systems for real industrial settings, taken from prototype to deployment.

The care-label QC system is a good example. I built it with OpenCV (template and feature matching, OCR, an explainable defect overlay) and iterated it three times, from a flatbed scanner to an industrial camera to a conveyor with a PLC reject. It cut inspection time by 99.5\% and improved throughput by up to 300\%. On the edge side, my published work at Trento runs real-time pattern detection in 14 ms on a Raspberry Pi 4, 74x faster than the DTW baseline, and I designed the benchmark suites and ablations that decided the architecture. I have also built the data and deployment pipelines around these models, using Python, PyTorch, Docker and AWS.

I should be direct about the gaps. My multimodal work is on audio, sensor and time-series data rather than image and text, so I have not yet fine-tuned vision-language models, and I have not worked with YOLO or Detectron2, or with Jetson, TensorRT and quantization in production. I have used HuggingFace Transformers and generative models (diffusion and VAE prototypes at McGill), and I expect to close these gaps quickly, since the evaluation, optimization and deployment fundamentals are the same.

Nexxa's focus on heavy industry is where I have done my best work, and I would like to bring that experience to your vision stack. I am a Canadian permanent resident based in \clCity{}, eligible to work without sponsorship, available immediately and comfortable working remotely. My portfolio, publications and open datasets are at \href{\clSite}{nishal.xyz}.

\vspace{12pt}
Sincerely,\\[2pt]
\textbf{\clName}, PhD

\end{document}
